\emph{Sistem bukti} mènèhi cara sintaktis kanggo mènèhi ciri
konsèkuènsi lan kompatibilitas antarané ukara. \emph{Dedhuksi
alami} iku salah sijiné. Ing sistem iki, \emph{derivasi} awujud
wit formula. Formula paling ndhuwur ing !!a{derivation} iku
\emph{asumsi}. Supaya !!{derivation} bener, formula liyané
kudu didhasaraké marang salah siji \emph{paugeran
inferensi}. Saben panggandhèng lan kuantor nduwèni paugeran
introduksi lan eliminasi. Upamané, yèn $!A$ diturunaké nganggo
paugeran $\Elim{\lif}$, formula sadurungé, yaiku \emph{prèmisé},
kudu $!B\lif !A$ lan~$!B$, kanggo sawijining~$!B$.
Sawetara paugeran bisa \emph{mbubaraké} asumsi. Upamané, yèn
$!A\lif !B$ dijupuk saka~$!B$ nganggo $\Intro{\lif}$,
kedadéan~$!A$ minangka asumsi ing !!{derivation} prèmis~$!B$
kena dibubaraké. Asumsi mau diwènèhi tandha kang uga dicathet
ing inferensi iku.

Yèn ana !!a{derivation} kanthi formula pungkasan~$!A$ lan
kabèh asumsi wis dibubaraké, $!A$ iku \emph{teorema},
ditulis $\Proves !A$. Yèn kabèh asumsi kang durung dibubaraké
ana ing himpunan~$\Gamma$, kita ngarani $!A$ \emph{bisa diderivasi
saka}~$\Gamma$ lan nulis $\Gamma\Proves !A$. Yèn
$\Gamma\Proves\lfalse$, $\Gamma$ \emph{ora konsistèn}.
Yèn ora, $\Gamma$ \emph{konsistèn}. Upamané, $\Gamma\Proves !A$
yèn lan mung yèn $\Gamma\cup\{\lnot !A\}$ ora konsistèn.
Sipat iki gegayutan karo semantik. Yèn $\Gamma\Proves !A$,
mula $\Gamma\Entails !A$. Iki diarani \emph{kesahihan}.
\emph{Teorema kesahihan} dibuktèkaké kanthi induksi tumrap
dawané !!{derivation}s. Saben inferensi njaga yèn kesimpulan
dadi konsèkuènsi saka asumsi kang isih kabukak, yèn prèmis-prèmisé
dadi konsèkuènsi saka asumsi-asumsi kang durung dibubaraké.
